\section{Open Questions}
\label{sec:open-questions}
The following five questions remain genuinely open after this replication. Each is
uncertainty-first: what we do \emph{not} know, why the replication does not settle it, and a
concrete probe that would.

\begin{enumerate}
\item \textbf{Ansatz generality: QAOA and UCCSD.} Do the paper's analytical-derivative formulas
produce correct forces for non-VQE variational ansatze (QAOA, UCCSD), where parameter-shift rules
require generalized (multi-frequency) shifts and the ansatz Hessian is more involved than for a
hardware-efficient $R_x/R_y$ + CNOT circuit?
\begin{itemize}
\item \emph{Basis:} Our replication used only the paper's 2-layer hardware-efficient ansatz. The
Hellmann--Feynman-at-optimum reduction is ansatz-agnostic in principle, but the second-order
response equations depend on ansatz structure; this is not tested.
\item \emph{Next steps:} Reimplement the H$_2$/STO-3G test in PennyLane with (a) $k$-UCCSD via
\texttt{qml.templates.UCCSD}, (b) a 2-layer QAOA-inspired ansatz. Compute analytical $dE/dR$ +
exact full-diag reference; report $|\Delta|$ vs our $2.9\times10^{-5}$ baseline and verify PennyLane
uses generalized parameter-shift for UCCSD.
\end{itemize}

\item \textbf{Shot / gate noise robustness.} How does the analytical VQE force degrade under
realistic shot noise and depolarizing gate noise, and at what shot budget does the analytical
approach lose its advantage over VQE-reoptimize numerical differences?
\begin{itemize}
\item \emph{Basis:} Our replication is noiseless (\texttt{default.qubit}). The paper's strongest
motivational statement is that analytical $\gg$ numerical \emph{under shot / gate noise}, which we
cannot test.
\item \emph{Next steps:} Repeat with \texttt{qml.device(\ldots, shots=N)} for $N\in\{10^2,10^3,10^4,10^5,10^6\}$,
20 seeds each. Plot $|\Delta_{\text{analytical}}|$ and $|\Delta_{\text{numerical}}|$ vs $N$. Repeat
on \texttt{default.mixed} with depolarizing 2Q error $p\in\{10^{-4},10^{-3},10^{-2}\}$. Extract
crossover $N$.
\end{itemize}

\item \textbf{End-to-end geometry optimization and vibrational analysis.} Does BFGS geometry
optimization driven by analytical VQE gradients (with optional analytical Hessian) converge to the
correct equilibrium for a triatomic like H$_2$O/STO-3G, and does the analytical Hessian give
correct harmonic frequencies?
\begin{itemize}
\item \emph{Basis:} This is the paper's primary downstream chemistry application (motivating VQE
forces at all), and is not exercised by our replication --- only first-derivative arithmetic at a
fixed geometry is verified.
\item \emph{Next steps:} (a) H$_2$ bond-length optimization from $r=1.0$~\AA{} via
\texttt{scipy.optimize.minimize(method="BFGS", jac=analytical\_vqe\_grad)}; expect convergence to
$r_{\text{eq}}\approx 0.735$~\AA{} within $5\times10^{-3}$~\AA{}. (b) H$_2$ harmonic frequency:
compute $d^2E/dR^2$ at $r_{\text{eq}}$ via response-equation solve; compare $\omega$ to experimental
$\sim4401$~cm$^{-1}$ and to PySCF analytical Hessian. (c) Extend to H$_2$O/STO-3G.
\end{itemize}

\item \textbf{Reaction-path robustness.} Does the analytical force stay accurate along a reaction
path (H + H$_2$ $\to$ H$_2$ + H, N$_2$ dissociation) where the VQE state may poorly approximate the
true eigenstate and the Hellmann--Feynman residual grows?
\begin{itemize}
\item \emph{Basis:} The HF simplification is exact only for exact eigenstates. At $r=0.735$~\AA{}
we have $|E_{\text{VQE}}-E_{\text{FCI}}|=5.75\times10^{-5}$~Ha (tiny). At stretched geometries in
larger systems, VQE can miss FCI by $\gg 10$~mHa, and the gradient error should grow correspondingly.
Untested.
\item \emph{Next steps:} Scan $r\in\{0.5,0.75,1.0,1.5,2.0,2.5\}$~\AA{} for H$_2$/STO-3G with
warm-started VQE; plot $|\Delta_{\text{analytical}}|$ vs $r$ and vs $|E_{\text{VQE}}-E_{\text{FCI}}|$.
Repeat for LiH/STO-3G (12 qubits) to probe multireference regime.
\end{itemize}

\item \textbf{Classical embedding.} Can the analytical VQE derivative framework be coupled to a
classical embedded environment (VQE-active-space in HF/DFT), and does the analytical force on the
QM atoms remain accurate when the environment's orbital response is included?
\begin{itemize}
\item \emph{Basis:} Real chemistry needs embedding. The paper derives for a pure-VQE $H(R)$;
embedding introduces environment-response terms that the paper does not address.
\item \emph{Next steps:} Build a CASCI(2e,2o)-in-HF embedding of H$_2$O in PennyLane+PySCF. Compare
(i) analytical HF on the active-space Hamiltonian only vs (ii) full PySCF CASCI analytical gradient
with environment response. Report the missing env-response contribution; propose an analog of the
CASSCF Z-vector correction term if needed.
\end{itemize}
\end{enumerate}
